\section{Application, Deployment, and Limitations}
The Gradio application presents criterion screening, evidence, an inference lab, benchmarks, and offline results. It is officially deployed as a locally hosted browser application, which satisfies the assignment. On Windows, \texttt{run\_app.bat} starts the app; the first run downloads the selected public model. An optional temporary Gradio share link is intended only for defense-day remote access. The Hugging Face Space remains staged but was not deployed because the platform requires a PRO account for the selected 2026 hardware.

The dataset contains a limited number of patients and the held-out split is correspondingly small; bootstrap intervals are included to communicate uncertainty. Labels and evidence inherit the source annotation's limitations, while evidence disagreement may reflect incomplete annotation rather than a model error. Synthetic-token inference benchmarks do not predict clinical workload latency. The application uses deterministic templates rather than independently verified rationales, and its aggregation rule is explicitly preliminary. CriterionLens is not a clinical decision system, must not be used for enrollment decisions, and requires expert review of every case.
